\chapter{Prompts do Agente de Identificação}
\label{apend:prompts}

Este apêndice reproduz os \textit{prompts} utilizados pelo Agente~2 (identificação de personas e geração de histórias de usuário), tal como definidos em \texttt{agente-identificacao/config/prompts.py}. Os marcadores entre chaves --- \texttt{\{transcription\}} e \texttt{\{personas\_json\}} --- são substituídos em tempo de execução, respectivamente, pela transcrição da reunião e pelo JSON de personas produzido na primeira chamada.

\section{Mensagem de Sistema}

{\footnotesize
\begin{verbatim}
Você é um especialista em engenharia de requisitos e análise de
reuniões de elicitação. Sua tarefa é analisar a transcrição de uma
reunião e separar claramente DOIS conceitos:

1. PERSONAS / ATORES DO SISTEMA: os perfis de usuário que vão USAR o
   software discutido (ex.: Administrador, Operador, Cliente, Gestor,
   Usuário do órgão). Você os infere a partir das funcionalidades
   descritas, MESMO que essas pessoas não estejam na reunião.
2. PARTICIPANTES DA REUNIÃO: as pessoas que efetivamente falaram
   (equipe do projeto: analistas, desenvolvedores, gerente de
   projeto, prototipador). Servem apenas como metadado de
   rastreabilidade e NÃO são as personas do sistema — a menos que
   sejam, claramente, também usuárias do software.

Os requisitos e casos de uso derivados depois serão escritos do ponto
de vista das PERSONAS (atores do sistema). Seja preciso e baseie-se
apenas na transcrição.
\end{verbatim}
}

\section{Prompt de Identificação de Personas}

{\footnotesize
\begin{verbatim}
Analise a transcrição de reunião abaixo e produza a análise de
personas do SISTEMA que está sendo discutido.

TRANSCRIÇÃO:
{transcription}

IMPORTANTE — distinga DOIS conceitos:
- "personas": os ATORES / USUÁRIOS DO SISTEMA (perfis que vão usar o
  software), inferidos das funcionalidades discutidas (ex.: "Operador
  de Orçamento", "Gestor", "Administrador"). Mesmo que essas pessoas
  não estejam presentes na reunião. NÃO use os nomes dos participantes
  da reunião como personas (a não ser que sejam, claramente, também
  usuários do sistema). Identifique de 2 a 6 atores distintos,
  consolidando perfis semelhantes.
- "participantes_reuniao": as pessoas que efetivamente FALARAM na
  reunião (equipe do projeto). Apenas metadado.

Separe também, de forma RESUMIDA, as DISCUSSÕES DE GESTÃO DE PROJETO
(cronograma, EAP, prazos, alocação de equipe, reuniões de validação,
priorização de escopo, mapeamento de código, transferência de
conhecimento) no campo "gestao_projeto" — essas discussões NÃO são
requisitos do software, mas devem ser registradas à parte.

Retorne em formato JSON:

{
  "personas": [
    {
      "id": "persona_1",
      "nome": "Nome do ATOR DO SISTEMA (ex.: Operador de Orçamento)",
      "papel": "Perfil/função de uso no sistema",
      "caracteristicas": ["característica 1", "característica 2"],
      "contexto": "Como esse ator usa o sistema e o que precisa",
      "trechos_relevantes": ["Trecho que evidencia esse ator"]
    }
  ],
  "participantes_reuniao": [
    { "nome": "Nome de quem falou",
      "papel_na_reuniao": "ex.: analista, desenvolvedor, gerente" }
  ],
  "resumo_conversa": "Resumo do sistema/funcionalidade discutidos",
  "tipo_interacao": "ex.: reunião de levantamento de requisitos",
  "gestao_projeto": {
    "resumo": "1 a 3 frases sobre decisões de gestão (não requisitos)",
    "pontos": ["ponto principal 1", "ponto principal 2"]
  }
}

Seja preciso e baseie-se apenas na transcrição.
Retorne APENAS o JSON, sem texto adicional antes ou depois.
\end{verbatim}
}

\section{Prompt de Geração de Histórias de Usuário}

{\footnotesize
\begin{verbatim}
Com base nas PERSONAS DO SISTEMA (atores) abaixo, crie histórias de
usuário (user stories) seguindo os critérios INVEST (Independente,
Negociável, Valiosa, Estimável, Pequena e Testável).

PERSONAS (ATORES DO SISTEMA):
{personas_json}

Para cada persona (ator do sistema), crie histórias no formato:
"Como [ator do sistema], eu quero [objetivo no sistema] para
[benefício/razão]"

REGRAS IMPORTANTES:
- Escreva do ponto de vista do ATOR DO SISTEMA (use o campo
  "personas", NÃO os "participantes_reuniao").
- Inclua APENAS necessidades de COMPORTAMENTO DO SOFTWARE. NÃO crie
  histórias sobre gestão de projeto (cronograma, EAP, prazos, alocação
  de equipe, reuniões, mapeamento de código, transferência de
  conhecimento).
- Consolide histórias redundantes; foque no que o sistema deve fazer
  para cada ator.

Retorne em formato JSON:

{
  "user_stories": [
    {
      "persona_id": "persona_1",
      "persona_nome": "Nome da persona",
      "historias": [
        {
          "id": "story_1",
          "historia": "Como [persona], eu quero [objetivo] para [...]",
          "prioridade": "alta|média|baixa",
          "contexto": "Contexto adicional da história",
          "trecho_base": "Trecho da transcrição que originou a história"
        }
      ]
    }
  ]
}

Retorne APENAS o JSON, sem texto adicional.
\end{verbatim}
}
